\subsection*{4.2 Rate amplification for dual-induced SLR}
In this section we describe our first rate amplification procedure for SLR that are induced by dual SLR. Unlike the previous section, it will be more convenient to explicitly state within the notation the set $P$ over which we are working as we will be dealing with several such sets. The field $\mathbb F$, however, remains suppressed from the notation as it remains fixed in all SLR under consideration. We start by defining the following map of functions.
\setcounter{nttheorem}{27}
\begin{ntdefinition}[Author Definition28]
Let $P$ be a set and $\mathbb F$ a field. For an integer $\ell\geq1$ we define the map $\Phi:(\mathbb F^P)^\ell\to\mathbb F^{P^\ell}$ as follows. Let $g_1,\ldots,g_\ell\in\mathbb F^P$. The function $\Phi(g_1,\ldots,g_\ell):P^\ell\to\mathbb F$ is defined by
\[\Phi(g_1,\ldots,g_\ell)(p_1,\ldots,p_\ell)=\prod_{i=1}^\ell g_i(p_i)\]
for every $(p_1,\ldots,p_\ell)\in P^\ell$.
\end{ntdefinition}
Observe that $\Phi$ is multi-linear. Further, when $\ell=2$ and $g_1,g_2$ are viewed as vectors rather than functions, $\Phi$ is the outer product of the vectors.
\begin{ntdefinition}[Author Definition29]
Let $P$ be a set, $\mathbb F$ a field. Let $L^P$ be a linear subspace of $\mathbb F^P$. For an integer $\ell\geq1$, we define
\[L^{P^\ell}=\operatorname{Span}\{\Phi(g_1,\ldots,g_\ell)\mid g_1,\ldots,g_\ell\in L^P\}.\]
\end{ntdefinition}
\begin{ntclaim}[Author Claim30]
With the notation of Definition29, $\dim(L^{P^\ell})\leq\dim(L^P)^\ell$.
\end{ntclaim}
\begin{proof}
Let $B=\{g_1,\ldots,g_b\}$ be a basis for $L^P$, where $b=\dim(L^P)$. Define
\[B'=\{\Phi(h_1,\ldots,h_\ell)\mid(h_1,\ldots,h_\ell)\in B^\ell\}.\]
Observe that to prove the claim, it suffices to show that for every $f_1,\ldots,f_\ell\in L^P$ it holds that $\Phi(f_1,\ldots,f_\ell)\in\operatorname{Span}(B')$. As $f_1,\ldots,f_\ell\in L^P$, for every $i\in[\ell]$ we can write $f_i=\sum_{j=1}^b\lambda_{i,j}g_j$ with $\lambda_{i,j}\in\mathbb F$. We have that
\begin{align*}
\Phi(f_1,\ldots,f_\ell)&=\Phi\left(\sum_{j_1=1}^b\lambda_{1,j_1}g_{j_1},\ldots,\sum_{j_\ell=1}^b\lambda_{\ell,j_\ell}g_{j_\ell}\right)\\
&=\sum_{j_1,\ldots,j_\ell\in[b]}\left(\prod_{t=1}^\ell\lambda_{t,j_t}\right)\cdot\Phi(g_{j_1},\ldots,g_{j_\ell}),
\end{align*}
where the last equality follows by the multi-linearity of $\Phi$.
\end{proof}
\begin{ntdefinition}[Author Definition31]
Let $P$ be a set, $\mathbb F$ a field, and let $(D^P,L^P)$ be $(q,\tau,\varepsilon,\rho)$-dual SLR. Let $\ell\geq1$ be an integer. For $p\in P^\ell$ we define the distribution $D_p^{P^\ell}$ as follows. Write $p=(p_1,\ldots,p_\ell)$. To sample an element from $D_p^{P^\ell}$ proceed as follows:
\begin{enumerate}
\item Sample $g_1\sim D_{p_1}^P,\ldots,g_\ell\sim D_{p_\ell}^P$ independently.
\item If there exists $i\in[\ell]$ such that $g_i=\bot$, return $\bot$; Otherwise
\item Return $\Phi(g_1,\ldots,g_\ell)$.
\end{enumerate}
The collection of distributions $\{D_p^{P^\ell}\mid p\in P^\ell\}$ is denoted by $D^{P^\ell}$.
\end{ntdefinition}
We have the following lemma.
\begin{ntlemma}[Author Lemma32]
Let $P$ be a set, $\mathbb F$ a field, and let $(D^P,L^P)$ be a $(q,\tau,\varepsilon,\rho)$-dual SLR. Let $\ell\geq1$ be an integer and $D^{P^\ell}$ as in Definition31. Then, for every $p,r\in P^\ell$,
\[\Pr_{g\sim D_p^{P^\ell}}[g(r)\ne0\mid g\ne\bot]\leq\tau^{\operatorname{dist}(p,r)}.\]
\end{ntlemma}
\begin{proof}
Write $p=(p_1,\ldots,p_\ell)$, $r=(r_1,\ldots,r_\ell)$. By Definition31, conditioned on $g\ne\bot$ we have that $g=\Phi(g_1,\ldots,g_\ell)$ with $g_i\sim D_{p_i}^P$ for each $i\in[\ell]$ independently. Thus, $g(r)\ne0$ is the event
\[\Phi(g_1,\ldots,g_\ell)(r_1,\ldots,r_\ell)=\prod_{i=1}^\ell g_i(r_i)\ne0.\]
By the independence of $g_1,\ldots,g_\ell$, and since we are working over a field $\mathbb F$ (and so a product is nonzero if and only if each of the terms is nonzero), we get
\begin{equation}\Pr_{g\sim D_p^{P^\ell}}[g(r)\ne0\mid g\ne\bot]=\prod_{i=1}^\ell\Pr_{g_i\sim D_{p_i}^P}[g_i(r_i)\ne0\mid g_i\ne\bot].\tag{4.1}\end{equation}
Let $T=\{i\in[\ell]\mid p_i\ne r_i\}$. As $D^P$ is a $(q,\tau,\varepsilon,\rho)$-dual SLR, for each $i\in T$ it holds that
\[\Pr_{g_i\sim D_{p_i}^P}[g_i(r_i)\ne0\mid g_i\ne\bot]\leq\tau.\]
Substituting to Equation(4.1), we get
\[\Pr_{g\sim D_p^{P^\ell}}[g(r)\ne0\mid g\ne\bot]\leq\tau^{|T|},\]
which completes the proof.
\end{proof}
\begin{ntdefinition}[Author Definition33]
Let $P$ be a set, $\mathbb F$ a field, and let $(D^P,L^P)$ be a $(q,\tau,\varepsilon,\rho)$-dual SLR. For an integer $\ell\geq1$ let $L^{P^\ell},D^{P^\ell}$ be as in Definition29 and Definition31, respectively. We denote the pair $(D^{P^\ell},L^{P^\ell})$ by $(D^P,L^P)^\ell$.
\end{ntdefinition}
\begin{ntproposition}[Author Proposition34]
Let $P$ be a set, $\mathbb F$ a field, and let $(D^P,L^P)$ be a $(q,\tau,\varepsilon,\rho)$-dual SLR. Then, for every integer $\ell\geq1$ we have that $(D^P,L^P)^\ell$ is a $(q_\ell,\tau_\ell,\varepsilon_\ell,\rho_\ell)$-dual SLR, where
\[q_\ell\leq q^\ell,\quad\tau_\ell\leq\tau,\quad\varepsilon_\ell\leq\ell\varepsilon,\quad\rho_\ell\geq1-(1-\rho)^\ell.\]
\end{ntproposition}
\begin{proof}
First note that for every $p\in P^\ell$, the distribution $D_p^{P^\ell}$ is supported on $F_p^{P^\ell}\cup\{\bot\}$. Indeed, if we write $p=(p_1,\ldots,p_\ell)$ then, conditioned on $g\ne\bot$, we have that $g=\Phi(g_1,\ldots,g_\ell)$ where $g_i\in D_{p_i}^P$. Thus,
\[g(p)=\Phi(g_1,\ldots,g_\ell)(p_1,\ldots,p_\ell)=\prod_{i=1}^\ell g_i(p_i)\ne0.\]
Moreover, by Definition29,
\[\bigcup_{p\in P^\ell}\operatorname{supp}(D_p^{P^\ell})\subseteq L^{P^\ell}\cup\{\bot\}.\]
We turn to show that $q_\ell\leq q^\ell$. Let $p=(p_1,\ldots,p_\ell)\in P^\ell$ and consider any $g\in\operatorname{supp}(D_p^{P^\ell})$. By Definition31, if $g\ne\bot$ then $g=\Phi(g_1,\ldots,g_\ell)$ where $g_i\in\operatorname{supp}(D_{p_i}^P)\setminus\{\bot\}$. Now, for every $r=(r_1,\ldots,r_\ell)\in P^\ell$ we have that
\[g(r)\ne0\quad\Longleftrightarrow\quad\prod_{i=1}^\ell g_i(r_i)\ne0.\]
Since $\mathbb F$ is a field, the above is equivalent to $g_i(r_i)\ne0$ for all $i\in[\ell]$. Hence there are at most $q^\ell$ points $r\in P^\ell$ for which $g(r)\ne0$, and so $q_\ell\leq q^\ell$.
The bound on the smoothness readily follows by Lemma32. Indeed, consider any pair of distinct $p,r\in\mathbb F^{P^\ell}$. We have that $\operatorname{dist}(p,r)\geq1$ and so, by Lemma32,
\begin{equation}\Pr_{g\sim D_p^{P^\ell}}[g(r)\ne0\mid g\ne\bot]\leq\tau^{\operatorname{dist}(p,r)}\leq\tau.\tag{4.2}\end{equation}
To bound the probability that $\bot$ is returned, note that the event $D^{P^\ell}=\bot$ holds only if for some $i\in[\ell]$, $g_i=\bot$. Hence, by the union bound, $\Pr[D_p^{P^\ell}=\bot]\leq\ell\varepsilon$. We conclude the proof by bounding the dimension of $L^{P^\ell}$. By assumption, $\dim(L^P)\leq(1-\rho)|P|$. Claim30 then implies that
\[\dim(L^{P^\ell})\leq\dim(L^P)^\ell\leq((1-\rho)|P|)^\ell=(1-\rho)^\ell|P^\ell|.\]
\end{proof}
\paragraph{Discussion on the smoothness $\tau_\ell=\tau$}
The downside of the rate amplification procedure that was given in this section is that $\tau_\ell$ does not decrease with $\ell$ (which is bad as, recall, we wish $\tau$ to be small as, by Claim22, the distance $\delta$ of the resulted LCC is proportional to $1/\tau$). Indeed, with the notation of Proposition34, $\tau_\ell=\tau$. By examining the proof and Lemma32 one natural idea is to consider an SLR not over the entire set $P^\ell$ but on some subset of it which is a code with distance, say, $d>1$. This will indeed guarantee that for every two points $p,r$ we have $\operatorname{dist}(p,r)\geq d$ and so the bound in Equation(4.2) will be $\tau^d$ rather than $\tau$. While natural, this idea fails to yield better parameters as the rate-loss incurred by using a code (even an MDS) is larger than the improvement on the rate guaranteed via the rate amplification procedure.
In the next sections we give a more elaborate rate amplification procedure (that is based on the one that was given in this section) in which $\tau$ does decrease with $\ell$. Roughly, $\tau_\ell=(q\cdot\log|P|)^{\operatorname{poly}(\ell)}\tau^\ell$, and so there is a slight loss in the smoothness, which the reader should think as negligible. The query complexity $q_\ell$ as well as the rate $\rho_\ell$ and $\varepsilon_\ell$ are all slightly worse than those obtained in the above rate amplification procedure and so the two techniques are incomparable.
